% 教学日序；方框标记交互，虚线是3日请求之前的信息边界。
\begin{tikzpicture}[
  x=1mm,y=1mm,font=\figurefont\small,
  event/.style={draw=BookInk,fill=BookBlue!5,line width=0.8pt,
    text width=21mm,minimum height=11mm,align=center,inner sep=1mm},
  future/.style={event,draw=BookAmber,dashed,fill=BookAmber!6},
  time/.style={-{Latex[length=2.2mm,width=1.5mm]},
    draw=BookMuted,line width=0.8pt},
  note/.style={font=\figurefont\footnotesize,text=BookMuted}
]
  \path[use as bounding box] (0,-2) rectangle (169,67);
  \foreach \x/\day in {30/1,60/2,90/3,120/4,150/5}{
    \node at (\x,53) {\day 日};
  }
  \node[anchor=west] at (0,35) {用户甲};
  \node[anchor=west] at (0,14) {用户乙};
  \draw[time] (20,35) -- (167,35);
  \draw[time] (20,14) -- (167,14);
  \draw[draw=BookTeal,dashed,line width=1pt] (76,0) -- (76,59);
  \node[anchor=south,text=BookTeal] at (76,61) {3日请求前的信息边界};
  \node[event] at (30,35) {1日交互\\可用历史};
  \node[event,fill=BookTeal!8,draw=BookTeal] at (90,35)
    {3日请求\\本次测试};
  \node[event] at (60,14) {2日交互\\可用历史};
  \node[future] at (120,14) {4日交互\\此时不可用};
  \node[future] at (150,14) {5日交互\\乙的测试};
  \node[note,anchor=north] at (135,5) {物品丙也到5日才上架};
\end{tikzpicture}
